%\documentclass[12pt,letterpaper]{article}



%%%
%% Required packages:
%%
\usepackage{etex}
\usepackage{amsmath}
\usepackage{amssymb}
\usepackage{amstext}
\usepackage{amsthm}
\usepackage{fancyhdr,fancybox}
\usepackage{mathrsfs} % need for \mathscr command
\usepackage{multicol}
\usepackage[normalem]{ulem}
%\usepackage{pictex,pictexplus}
% \usepackage{pictexwd}
\usepackage{rotating}
\usepackage{url}
\usepackage{xfrac}
\usepackage{booktabs}
\usepackage{color,soul}
\usepackage{comment}

\usepackage{hyperref}
\hypersetup{
    colorlinks=true,     % false: boxed links; true: colored links
    linkcolor=blue,      % color of internal links
    citecolor=blue,      % color of links to bibliography
    filecolor=blue,      % color of file links
    urlcolor=blue        % color of external links
}


\usepackage{tikz}
\usetikzlibrary{backgrounds,calc}

%%
%% Common formatting commands
%%
\setlength{\parindent}{0in}
\setlength{\textwidth}{7in}
\setlength{\evensidemargin}{-0.25in}
\setlength{\oddsidemargin}{-0.25in}
\setlength{\parskip}{.5\baselineskip}
\setlength{\topmargin}{-0.5in}
\setlength{\textheight}{9in}

%               Problem and Part
\newcounter{problemnumber}
\newcounter{partnumber}
\newcounter{subpartnumber}
\newcommand{\Problem}{%
  \refstepcounter{problemnumber}%
  \setcounter{partnumber}{0}%
  \setcounter{subpartnumber}{0}%
  \item[\makebox{\hfil\textbf{\theproblemnumber.}\hfil}]%
  }
\newcommand{\InProblem}[2][1.6]{%
  \refstepcounter{problemnumber}%
  \setcounter{partnumber}{0}%
  \setcounter{subpartnumber}{0}%
  \textbf{\theproblemnumber.}\ \ \parbox[t]{#1in}{#2}%
}
\newcommand{\InSmallProblem}[1]{\InProblem[1.05]{#1}}
\newcommand{\NoProblem}[1]{%
  \mbox{\hfil\textbf{#1}\hfil}%
  }
\newcommand{\UnProblem}{%
  \refstepcounter{problemnumber}%
  \setcounter{partnumber}{0}%
  \setcounter{subpartnumber}{0}%
  \mbox{\hfil\textbf{\theproblemnumber.}\hfil}%
  }
\newcommand{\InFlexProblem}[1]{%
  \refstepcounter{problemnumber}%
  \setcounter{partnumber}{0}%
  \setcounter{subpartnumber}{0}%
  \textbf{\theproblemnumber.}\ \ #1%
}
%%  Part:
\newcommand{\Part}{%
  \renewcommand{\thepartnumber}{(\alph{partnumber})}%
  \refstepcounter{partnumber}
  \setcounter{subpartnumber}{0}%
  \item[(\alph{partnumber})]%
}
\newcommand{\InPart}[2][1.75]{%
  \renewcommand{\thepartnumber}{(\alph{partnumber})}%
  \refstepcounter{partnumber}%
  \setcounter{subpartnumber}{0}%
  (\alph{partnumber})\ \ \parbox[t]{#1in}{#2}%
}
\newcommand{\UnPart}{%
  \refstepcounter{partnumber}%
  \renewcommand{\thepartnumber}{(\alph{partnumber})}%
  \setcounter{subpartnumber}{0}%
  (\alph{partnumber})%
}
\newcommand{\InLargePart}[1]{\InPart[3.05]{#1}}
\newcommand{\InFullPart}[1]{\InPart[6.2]{#1}}
\newcommand{\InSmallPart}[1]{\InPart[1.05]{#1}}
%% SubPart:
\newcommand{\SubPart}{%
  \refstepcounter{subpartnumber}%
  \item[(\roman{subpartnumber})]%
  }
\newcommand{\InSubPart}[2][2.25]{%
  \stepcounter{subpartnumber}%
  (\roman{subpartnumber})\ \ \parbox[t]{#1in}{#2}%
}
\newcommand{\InSmallSubPart}[1]{\InSubPart[1.5]{#1}}
\newcommand{\InTinySubPart}[1]{%
  \stepcounter{subpartnumber}%
  (\roman{subpartnumber})\ \ \makebox{#1}%
}
\newcommand{\InMySubPart}[2][2.25]{%
  \stepcounter{subpartnumber}%
  \parbox[t]{#1in}{\makebox[0.3in][l]{(\roman{subpartnumber})}\ \ {#2}}%
}
\newcommand{\TrueFalse}{\textbf{T}\ \ \textbf{F}\ \ }
\newcommand{\TFPart}[1]{% 
  \stepcounter{partnumber}%
  \setcounter{subpartnumber}{0}%
  \item[(\alph{partnumber})] \TrueFalse\parbox[t]{5.5in}{#1}}

\newcommand{\Points}[1]{(#1 points) \ }
\newcommand{\Point}[1]{(#1 point) \ }


%% For Answers:
\newcounter{answernumber}
\newcommand{\Answer}{%
  \refstepcounter{answernumber}%
  \setcounter{partnumber}{0}%
  \setcounter{subpartnumber}{0}%
  \item[\makebox{\hfil\textbf{\theanswernumber.}\hfil}]%
  }


% Setting up the headers for the first page
%\chead{} % will default to what is typed in the file.
\fancyhead[L]{\textsc{Math 22a}}
\fancyhead[R]{\textsc{Fall 2024}}
\fancyfoot[R]{
	\vspace*{\fill}\footnotesize\textcopyright{} \emph{Harvard Math 22a}	
}
\renewcommand{\headrulewidth}{0pt}

%\fancyhead[C]{}
%	\chead{}	
%	\lhead{}
%	\rhead{}
%	\cfoot{}


% Putting copyright on the footer of each page
%\fancypagestyle{secondpage}{
%	\fancyhead[C]{TESTing again} % change this to be blank
%		%\chead{}	
%	\fancyhead[L]{bears} % change this to be blank
%	\fancyhead[R]{dogs} % change this to be blank
%	\fancyfoot[R]{ %keeps the footer the same
%		\vspace*{\fill}\footnotesize\textcopyright{} \emph{Harvard Math 22a}	
%	}
%}
%\renewcommand{\headrulewidth}{0pt}
%\pagestyle{fancy}
\pagestyle{empty}


%\lhead{}
%\rhead{}
%\rfoot{\vspace*{\fill}\footnotesize\textcopyright{} \emph{Harvard Math 22a}}
%\renewcommand{\headrulewidth}{0pt}
%\pagestyle{fancy}




%%
%% Common shortcut commands:
%%
\newcommand{\ds}{\displaystyle}
\newcommand{\sgn}{\operatorname{sgn}}
\renewcommand{\Pr}{\operatorname{P}}
\newcommand{\CPr}[3][{}]{\Pr_{\text{#1}}\left(#2\,|\,#3\right)}

\newcommand{\C}{\mathbb{C}}
\newcommand{\R}{\mathbb{R}}
\newcommand{\Z}{\mathbb{Z}}
\newcommand{\N}{\mathbb{N}}
\newcommand{\Q}{\mathbb{Q}}
\newcommand{\D}{\mathbb{D}}

\newcommand{\rref}{\operatorname{rref}}
\newcommand{\im}{\operatorname{im}}
\newcommand{\rank}{\operatorname{rank}}
\newcommand{\diag}{\operatorname{diag}}
\newcommand{\Span}{\operatorname{span}}
%\renewcommand{\vec}[1]{\mathbf{#1}}
\newcommand{\charpoly}{\mathscr{P}}
\newcommand{\nicefrac}[2]{\sfrac{#1}{#2}} % using xfrac package
\newcommand{\topology}{\mathcal{T}}
\newcommand{\Inn}{\operatorname{Inn}}

\newcommand{\blankbox}[2]{\fbox{\rule{#1}{0in}\rule{0in}{#2}}}

\newenvironment{myindentpar}[1]%
 {\begin{list}{}%
         {\setlength{\leftmargin}{#1}
          \setlength{\rightmargin}{\leftmargin}}%
         \item[]%
 }
 {\end{list}}
\newtheorem{theorem}{Theorem}
\newtheorem*{NStheorem}{Theorem}
\newtheorem{cor}[theorem]{Corollary}
\newtheorem*{claim}{Claim}
\newtheorem{defn}{Definition}

\newenvironment{amatrix}[1]{%
  \left(\begin{array}{@{}*{#1}{c}|c@{}}
}{%
  \end{array}\right)
}

%--------grstep
% For denoting a Gauss' reduction step.
% Use as: \grstep{\rho_1+\rho_3} or \grstep[2\rho_5 \\ 3\rho_6]{\rho_1+\rho_3}
\newcommand{\grstep}[2][\relax]{%
   \ensuremath{\mathrel{
       {\mathop{\longrightarrow}\limits^{#2\mathstrut}_{
                                     \begin{subarray}{l} #1 \end{subarray}}}}}}
\newcommand{\swap}{\leftrightarrow}




\section{{Homogeneous Systems, $\vec\S$1.5}}% 
\chead{\Large\textbf{Homogeneous Systems, $\vec\S$1.5}}% 

%\begin{document}
\thispagestyle{fancy}


Recall from last class:

\begin{itemize}

\item Let $A$ be an $m \times n$ matrix, then $$A = [ \vec{a}_1\; \dots \; \vec{a}_n], \; \; \vec{a}_j \in \mathbb{R}^m,\; 1\leq j \leq n.$$  
If $x \in \mathbb{R}^n$, then we define $$A \vec{x} = \vec{a}_1 x_1 + \vec{a}_2 x_2 +\dots + \vec{a}_n x_n.$$ Thus the matrix equation $A\vec{x} =\vec{b}$ is the same as the vector equation $$\vec{a}_1 x_1 + \vec{a}_2 x_2 +\dots + \vec{a}_n x_n=\vec{b}.$$
\item $A\vec{x} =\vec{b}$ has a solution if and only if $\vec{b} \in \text{Span}(\vec{a}_1,\dots, \vec{a}_n)$ where $\vec{a}_j$ is the $j$-th column of $A$.
\item A set of vectors $\vec{v}_1, \dots, \vec{v}_p$ in $\mathbb{R}^m$ {\bf spans} $\mathbb{R}^m$ if every vector in $\mathbb{R}^m$ is a linear combination of $\vec{v}_1, \dots, \vec{v}_p$.
\item {\bf Theorem (Theorem 4 in the book)} Let $A$ be an $m \times n$ matrix, then the following are equivalent.
\begin{enumerate}
\item For each $\vec{b} \in \mathbb{R}^m$, the matrix equation $A \vec{x} =\vec{b}$ is consistent.
\item Each $\vec{b} \in \mathbb{R}^m$ is a linear combination of the columns of $A$.
\item The columns of $A$ span $\mathbb{R}^m$.
\item $A$ has a pivot position in every row.
\end{enumerate}
\item {\bf Fact (or Theorem 5 from the book):} If $A$ is a $m \times n$ matrix and $\vec{u}, \vec{v} \in \mathbb{R}^n$, $c\in \mathbb{R}$, then
\begin{enumerate}
\item $A(\vec{u} +\vec{v})=A\vec{u} +A\vec{v}$
\item $A(c \vec{u}) =c (A\vec{u}).$
\end{enumerate}
\end{itemize}

\newpage


\begin{defn}
A linear system is called {\bf homogeneous} if it can be written as $A\vec{x}=\vec{0}$ where $A$ is an $m \times n$ matrix, $\vec{x} \in \mathbb{R}^n$, and $\vec{0}$ is the vector of $m$ zeros. 
\end{defn}

\begin{enumerate}

\Problem  Let $A=\begin{bmatrix} 2 & -5 & 8\\ -2 & -7 & 1\\ 4& 2 & 7\\ \end{bmatrix}$. Find all solutions to $A\vec{x} = \vec{0}$.

\newpage

{\bf Theorem:} If $A\vec{x} =\vec{b}$ is consistent for a given $\vec{b}$, and $\vec{p}$ is a solution, then the solution set of $A\vec{x}=\vec{b}$ is the set of vectors of the form $\vec{p}+\vec{v}_h$ where $\vec{v}_h$ is a solution of $A\vec{x} =\vec{0}$.

\Problem Prove the Theorem.

\newpage

\Problem Let $A=\begin{bmatrix} 1 & 0 & 5\\ -2 & 1 & -6\\ 0 & 2 & 8\\ \end{bmatrix}$. 
\Part Find all solutions to $A\vec{x} = \vec{0}$.

\vfill

\Part Let $\vec{b}=\begin{bmatrix} 2\\-1\\6\\\end{bmatrix}$. Describe all solutions to $A\vec{x} =\vec{b}$.

\vfill

\Problem Construct a 2 by 2 matrix $A$ such that the solution set to $A\vec{x}=\vec{0}$ is the line in $\mathbb{R}^2$ going through $(0,0)$ and $(4,1)$.

\vfill


\newpage

\begin{defn}
A set of vectors $\{ \vec{v}_1, \dots, \vec{v}_p\}$ in $\mathbb{R}^n$ is said to be {\bf linearly independent} if $$x_1 \vec{v}_1 +x_2 \vec{v}_2 +\dots +x_p \vec{v}_p = \vec{0}$$ has only the trivial solution. The set is called {\bf linearly dependent} if the equation has non-trivial solutions and the equation is called the {\bf dependence relation}.
\end{defn}

\Problem Let $\vec{v}_1=\begin{bmatrix} 1\\2\\3\\ \end{bmatrix}$, $\vec{v}_2=\begin{bmatrix} 4\\5\\6\\ \end{bmatrix}$, and $\vec{v}_3=\begin{bmatrix} 2\\1\\0\\ \end{bmatrix}$. Is $\{ \vec{v}_1, \vec{v}_2,\vec{v}_3 \}$ linearly independent?

\vfill

\vfill

\vfill


Let $A= [\vec{a}_1 \; \dots \; \vec{a}_n ]$, then $\{ \vec{a}_1 \; \dots \; \vec{a}_n \}$ are linearly independent if and only if $A \vec{x} = \vec{0}$ has only the trivial solution.

\Problem Is a set with only one vector always linearly independent?

\vfill

\Problem If $\vec{0} \in \{ \vec{v}_1, \dots, \vec{v}_p \}$, then is the set linearly dependent?

\vfill

\Problem If $\vec{0} \in \text{Span}\{ \vec{v}_1, \dots, \vec{v}_p \}$, then the vectors are linearly independent?

\vfill

%End here for worksheet without sols.
\end{enumerate}
%\end{document}
%%


% This was moved to the next worksheet
\begin{comment} 
\newpage

\Problem Let $\vec{v}_1=\begin{bmatrix} 1\\2\\ \end{bmatrix}$, $\vec{v}_2=\begin{bmatrix} 3\\4\\ \end{bmatrix}$, and $\vec{v}_3=\begin{bmatrix} 5\\6\\ \end{bmatrix}$. Are the vectors linearly independent?

\vfill

{\bf Theorem 7:} A set $S=\{ \vec{v}_1, \dots, \vec{v}_p \}$ of two or more vectors is linearly dependent if and only if one of the vectors of $S$ is a linearly combination of the others.

Moreover if $\vec{v}_1\neq \vec{0}$, then some $\vec{v}_j$ for $j>1$ is a linear combination of $\vec{v}_1, \dots, \vec{v}_{j-1}$.

\Problem Prove Theorem 7.

\vfill

\newpage

\Problem Let $\vec{u}=\begin{bmatrix} 1\\ 0\\ 1\\ \end{bmatrix}$ and $\vec{v} = \begin{bmatrix} 2 \\ 0 \\ 8\\ \end{bmatrix}$
\Part Describe $\text{Span}(\vec{u}, \vec{v})$

\vfill

\Part Prove $\vec{w} \in Span(\vec{u}, \vec{v})$ if and only if $\{ \vec{u}, \vec{v}, \vec{w} \}$ is linearly dependent.

\vfill
\vfill

{\bf Theorem 8:} Any set $\{ \vec{v}_1, \dots, \vec{v}_p \}$ in $\mathbb{R}^n$ is linearly dependent if $p>n$.

\Problem Prove the theorem.

\vfill
\vfill
\vfill 
\end{comment} 
% the above was moved to a different worksheet.


%\end{enumerate}


%\end{document}

\begin{comment}
\newpage
\chead{\Large\textbf{Homogeneous Systems and Linear Independence -- Answers / Solutions}}%
\lhead{}
\rhead{}
\thispagestyle{fancy}

\begin{enumerate}
  \Answer The REF for $[A|0]$ (here $0$ means the $3\times 1$ column vector $0\in \mathbb{R}^3$) is:
  \[
    \left[
      \begin{array}{ccc|c}
        1&0&\frac{17}{8}&0\\
        0&1&-\frac{3}{4}&0\\
        0&0&0&0
      \end{array}
    \right]
  \] There are two basic variables $x_1,x_2$ and one free variable $x_3$. The general solution is
  \[
    x_3\left[
      \begin{array}{c}
        -\frac{17}{8}\\
        \frac{3}{4}\\
        1
      \end{array}
    \right]
  \]

  \Answer Consider the two sets $S_1=\{\vec{x}\mid A\vec{x}=\vec{b}\}$ and $S_2=\{\vec{p}+\vec{v}_h\mid A\vec{v}_h=0\}$. The first set $S_1$ is just the usual solution set of $A\vec{x}=\vec{b}$, and the second set $S_2$ is the set described in the problem. Our goal is to show that these two things are equal. The way we will do this is by first showing that $S_1\subset S_2$ (that every element of $S_1$ is contained in $S_2$) and $S_2\subset S_1$ (every element of $S_2$ is contained in $S_1$).

  $S_1\subset S_2$: Let $\vec{x}\in S_1$. We also know that $A\vec{p}=\vec{b}$. Therefore, $A(\vec{x}-\vec{p})=A\vec{x}-A\vec{p}=\vec{b}-\vec{b}=0$. So let $\vec{v}_h=\vec{x}-\vec{p}$. Then $\vec{p}+\vec{v}_h=\vec{p}+\vec{x}-\vec{p}=\vec{x}$, so $\vec{x}\in S_2$.

  $S_2\subset S_1$: Let $\vec{p}+\vec{v}_h\in S_2$. Apply $A$ to this element: $A(\vec{p}+\vec{v}_h)=A\vec{p}+A\vec{v}_h=\vec{b}+\vec{0}=\vec{b}$, so we have that $\vec{p}+\vec{v}_h\in S_1$ is a solution to $A\vec{x}=\vec{b}$.

  Therefore, since both $S_1\subset S_2$ and $S_2\subset S_1$, then they must be equal.

  \Answer
  \begin{enumerate}
  \item The REF of the augmented matrix $[A|\vec{0}]$ is:
    \[
      \left[
        \begin{array}{ccc|c}
          1&0&5&0\\
          0&1&4&0\\
          0&0&0&0
        \end{array}
      \right]
    \]
    There are two basic variables $x_1,x_2$ and one free variable $x_3$. The general solution is:
    \[
      x_3\left[
        \begin{array}{c}
          -5\\
          -4\\
          1
        \end{array}
      \right]
    \]
  \item The REF of the augmented matrix $[A|\vec{b}]$ is:
    \[
      \left[
        \begin{array}{ccc|c}
          1&0&5&2\\
          0&1&4&3\\
          0&0&0&0
        \end{array}
      \right]
    \]
    Again, there are two basic variables $x_1,x_2$ and one free variable $x_3$. The general solution is:
    \[
      \left[
        \begin{array}{c}
          2\\
          3\\
          0
        \end{array}
      \right]+
      x_3\left[
        \begin{array}{c}
          -5\\
          -4\\
          1
        \end{array}
      \right]
    \]
  \end{enumerate}


    \Answer
    The line passing through $(0,0)$ and $(4,1)$ has equation $y=\frac{1}{4}x$. We can also write this as $x-4y=0$. Therefore, the matrix
    \[
      \left[
        \begin{array}{cc}
          1&-4\\
          0&0
        \end{array}
      \right]
    \] has the correct solution set. 


\Answer No, the three vectors are not linearly independent.  
 \begin{enumerate}
 \item The problem of linear independence is equivalent to testing whether the matrix equation $[\vec v_1\  \vec v_2\  \vec v_3] \vec x = \vec 0$, which has corresponding augmented matrix:
    \[
      \left[
        \begin{array}{ccc|c}
          1&4&2&0\\
          2&5&1&0\\
          3&6&0&0
        \end{array}
      \right]
    \]
 
  \item The REF of the augmented matrix $[\vec v_1\  \vec v_2\  \vec v_3|\vec{0}]$ is:
    \[
      \left[
        \begin{array}{ccc|c}
          1&0&-2&0\\
          0&1&1&0\\
          0&0&0&0
        \end{array}
      \right]
    \]
    There are two basic variables $x_1,x_2$ and one free variable $x_3$, meaning that there are infinitely many solutions, which means that there exists a non-trivial (i.e., not the zero vector) solution.
    
    \item If there is one non-trivial solution, the vectors are not linearly dependent, and 
   \[
   \left[
        \begin{array}{c}
          2\\
          -1\\
          1
        \end{array}
      \right]
   \]
   gives one such solution, which corresponds to the fact that $2 \vec v_1 - \vec{v}_2 + \vec v_3 = \vec 0$.
\end{enumerate}

\begin{comment}

\Answer No, the set containing only the zero vector, $\{ \vec 0\}$ is not linearly independent, since $22(\vec 0) = \vec 0$.  However, it is true that a set containing only one \emph{non-zero} vector is linearly independent.

\Answer No, such a set would not be linearly independent.  Say, for example, that $v_j = \vec 0$.  Then $0 \vec v_1 +\dots + 0 \vec v_{j-1} + 22 \vec v_j + 0 \vec v_{j+1}+\dots+0 \vec v_p = \vec 0$, showing that the set is not linearly independent.

\Answer The set of vector may or may not be linearly independent.  Not that $\vec 0 \in 
\text{Span} S$ for any set of vectors $S$ (just multiply all the vectors by zero).  So the hypothesis here is vacuous, and the set $S$ may or may not be linearly independent.

%\end{comment}


\end{enumerate}

\end{comment}
%\end{document}


%%% Local ables: 
%%% mode: latex
%%% TeX-master: t
%%% End: 
